\begin{figure}[t]
\centering
\begin{tikzpicture}[
  node distance=2.5cm and 3.5cm,
  every node/.style={font=\sffamily},
  box/.style={
    rectangle,
    draw=black,
    line width=1.2pt,
    minimum width=2.4cm,
    minimum height=1.2cm,
    align=center,
    text width=2.2cm
  },
  optbox/.style={
    rectangle,
    draw=black,
    line width=1pt,
    minimum width=2.2cm,
    minimum height=0.7cm,
    align=center,
    text width=2cm,
    font=\small
  },
  perfbox/.style={
    rectangle,
    draw=black,
    line width=1pt,
    minimum width=2.0cm,
    minimum height=0.55cm,
    align=center,
    text width=1.8cm,
    font=\tiny
  },
  arrow/.style={
    ->,
    line width=1.2pt
  },
  lbl/.style={
    font=\small,
    align=center
  }
]

% Stage 1: PyTorch Model
\node[box, fill=Garnet!20, draw=Garnet] (pytorch) at (0,0) {
  \textbf{PyTorch Model}\\
  (FP32)
};

% Performance annotation for PyTorch
\node[below=0.3cm of pytorch, lbl] (pytorch_perf) {
  \textcolor{Congaree}{\textbf{8--12 FPS}}
};

% Arrow with label: Export
\draw[arrow, black!90] (pytorch.east) -- ++(1.4,0) node[midway, above, lbl] {Export};

% Stage 2: ONNX
\node[box, fill=Garnet!20, draw=Garnet] at (4.8,0) (onnx) {
  \textbf{ONNX Graph}
};

% Arrow with label: TensorRT Compile
\draw[arrow, black!90] (onnx.east) -- ++(1.4,0) node[midway, above, lbl] {TRT Compile};

% Stage 3: TensorRT Engine
\node[box, fill=Garnet!20, draw=Garnet] at (8.8,0) (tensorrt) {
  \textbf{TensorRT Engine}\\
  (INT8)
};

% Performance annotation for final engine
\node[below=0.3cm of tensorrt, lbl] (engine_perf) {
  \textcolor{Congaree}{\textbf{95--125 FPS}}\\
  \textcolor{Congaree}{\textbf{(10--15$\times$)}}
};

% Title for optimization stages
\node[below=2.5cm of onnx, lbl, font=\normalsize] (opt_title) {
  \textbf{TensorRT Optimization Stages:}
};

% Optimization substages
\node[optbox, below=0.5cm of opt_title, fill=Atlantic!15, draw=Atlantic] (opt1) {
  1. Operator Fusion
};
\node[optbox, below=0.15cm of opt1, fill=Atlantic!15, draw=Atlantic] (opt2) {
  Conv + BN + ReLU $\to$ 1 kernel
};

\node[optbox, below=0.4cm of opt2, fill=Atlantic!15, draw=Atlantic] (opt3) {
  2. Constant Folding
};
\node[optbox, below=0.15cm of opt3, fill=Atlantic!15, draw=Atlantic] (opt4) {
  Pre-compute static ops
};

\node[optbox, below=0.4cm of opt4, fill=Atlantic!15, draw=Atlantic] (opt5) {
  3. Quantization
};
\node[optbox, below=0.15cm of opt5, fill=Atlantic!15, draw=Atlantic] (opt6) {
  FP32 $\to$ INT8 + calib.
};

\node[optbox, below=0.4cm of opt6, fill=Atlantic!15, draw=Atlantic] (opt7) {
  4. Kernel Auto-Tuning
};
\node[optbox, below=0.15cm of opt7, fill=Atlantic!15, draw=Atlantic] (opt8) {
  Benchmark on target HW
};

% Title for performance milestones
\node[below=1.0cm of opt8, lbl, font=\normalsize] (milestone_title) {
  \textbf{Performance Milestones:}
};

\node[perfbox, below=0.25cm of milestone_title, fill=Congaree!15, draw=Congaree] (mile1) {
  Graph Opt: 22--28 FPS (2.8$\times$)
};

\node[perfbox, below=0.15cm of mile1, fill=Congaree!15, draw=Congaree] (mile2) {
  INT8 Quant: 95--125 FPS (4.3$\times$)
};

\node[perfbox, below=0.15cm of mile2, fill=Congaree!15, draw=Congaree] (mile3) {
  Batching: 1500--2000 FPS TP
};

% Vertical dashed lines connecting stages to optimization box
\draw[dashed, black!90, line width=0.8pt] ($(onnx.south) + (0, -0.5)$) -- ($(opt_title.north) + (0, 0.2)$);
\draw[dashed, black!90, line width=0.8pt] ($(tensorrt.south) + (0, -0.5)$) -- ($(opt_title.north) + (0, 0.2)$);

\end{tikzpicture}

\caption{TensorRT optimization pipeline for edge deployment. A PyTorch-trained YOLOv8 model is exported to ONNX, then compiled by TensorRT with operator fusion, constant folding, INT8 quantization, and hardware-specific kernel tuning, achieving 10--15$\times$ overall speedup on Jetson platforms.}
\label{fig:tensorrt_pipeline_wrapped}
\end{figure}
